\section{A proper-volume obstruction}\label{sec:proper-volume}

The obstruction uses the fact that the original volume form extends
regularly to a proper model over the line of phase values. The
\(V\)-filtration records this regularity at the zero fiber. A nonzero
polynomial period would force the same volume class to have a nonzero
image in a module supported at zero, where its filtration degree is
impossible.

\begin{lemma}\label{lem:proper-volume}
Let \(Y\) be a smooth irreducible complex algebraic variety of dimension
\(n\), let \(f:Y\to\mathbb A^1_s\) be a proper nonconstant morphism, and
let \(\omega\) be a regular algebraic \(n\)-form on \(Y\). For any compact
closed \(n\)-cycle \(\gamma\) on \(Y(\mathbb C)\), the entire function
\[
 P_\gamma(\tau)=\int_\gamma e^{-\tau f}\omega
\]
cannot be a nonzero polynomial.
\end{lemma}

\begin{proof}
Factor \(f=\pi\circ i_f\), with \(i_f:Y\hookrightarrow
Y\times\mathbb A^1_s\) its graph and \(\pi\) the projection. Put
\(\mathcal B_f=i_{f+}\mathcal O_Y\) and
\(\mathcal M=\mathcal H^0(\pi_+\mathcal B_f)\), using left
\(\mathcal D\)-modules. Use the decreasing Kashiwara--Malgrange
filtration along \(s=0\), with
\(\partial_s s-\alpha\) nilpotent on \(\operatorname{gr}_V^\alpha\).
The graph generator \(e_f=\delta(s-f)\) belongs to
\(V^{>0}\mathcal B_f\). This is the graph-generator consequence of
strict negativity of the Bernstein--Sato roots
\cite[Corollary~5.2]{kashiwara-b-functions}
\cite[Corollary~5.15 and Example~5.19]{chen-dirks-mustata}.
In particular at each point of the graph over zero the generator
lies in \(V^\alpha\) for some rational \(\alpha>0\) locally.
Quasi-compactness of the fiber gives a common positive index
by taking the minimum over finitely many neighborhoods.
Off \(s=0\) every filtration piece is the whole graph module, and
off the graph there is nothing to check. Thus the top relative
form \(\omega e_f\) represents a degree-zero class in the relative
de Rham complex with coefficients in \(V^\alpha\mathcal B_f\).
The direct-image theorem for a smooth projection
\cite[Theorem~8.1]{chen-dirks-mustata} requires properness on the support
of the module. Here the support is the graph and the restriction of
\(\pi\) to it identifies with the proper map \(f\). The theorem therefore
identifies the image of the filtered complex with the corresponding
filtration piece of \(\mathcal M\). Consequently the class
\(m=[\omega e_f]\) satisfies
\[
 m\in V^\alpha\mathcal M\quad\text{at }s=0,\qquad \alpha>0.
\]
This step applies to the actual volume class even if the zero fiber
has a positive-dimensional critical locus.

We next check that the period respects the relations on the Fourier
transform of \(m\). Write \(M=\Gamma(\mathbb A^1_s,\mathcal M)\).
The graph action \cite[Equation~(27)]{chen-dirks-mustata} identifies
the graph module along \(Y\) with \(\mathcal O_Y[\partial_s]e_f\),
with actions
\[
 \begin{split}
 v(a\partial_s^j e_f)&=v(a)\partial_s^j e_f
                     -a\,v(f)\partial_s^{j+1}e_f,\\
 s(a\partial_s^j e_f)&=fa\partial_s^j e_f
                     -j a\partial_s^{j-1}e_f
 \end{split}
\]
for any local vector field \(v\) on \(Y\).
The Fourier convention is \(\tau=\partial_s\) and
\(\partial_\tau=-s\), which preserves the Weyl relation. Therefore
\(\widehat M\) is represented by degree-\(n\) hypercohomology on \(Y\) of
\[
 \bigl(\Omega_Y^\bullet[\tau],\,d_Y-\tau\,df\wedge\bigr),
 \qquad \partial_\tau\text{ acting by }\partial/\partial\tau-f.
\]
Indeed the relative de Rham complex of the projection uses the shift
by \(n\), and global sections on the affine target are exact.
One can compute this hypercohomology on a finite affine open cover of
\(Y\), whose finite intersections are affine since \(Y\) is separated.
This description uses the full polynomial parameter ring
\(\mathbb C[\tau]\).

At any \(\tau_0\in\mathbb C\), substituting \(\tau=\tau_0\)
in the total \v{C}ech complex commutes with both differentials.
After analytification, multiplying by \(e^{-\tau_0 f}\) turns the
specialized twisted differential into the ordinary holomorphic de
Rham differential. The affine intersections analytify to Stein
spaces, so the resulting \v{C}ech complex computes holomorphic de
Rham hypercohomology. Its class pairs with \(\gamma\) by analytic
de Rham comparison. A \v{C}ech coboundary thus has zero period
even when there is no global algebraic primitive. For a Weyl
operator \(A\), the class \(A\widehat m\) is represented by a global
polynomial family of top forms. Differentiating its exponential
integral on the fixed compact cycle realizes \(\partial_\tau\)
as \(\partial/\partial\tau-f\) on such a family. It follows that
evaluation on \(\gamma\) sends \(A\widehat m\) to \(A P_\gamma\)
and respects all relations on the cyclic submodule
\(\widehat N=\mathbb C[\tau]\langle\partial_\tau\rangle\,\widehat m\).

If \(P_\gamma\) were a nonzero polynomial, it would generate the
polynomial module \(\mathbb C[\tau]\) under Weyl operators, since
differentiation first produces a nonzero constant. The period
evaluation would give a surjection
\(\widehat N\to\mathbb C[\tau]\) taking \(\widehat m\) to \(P_\gamma\).
Inverse Fourier transform sends the polynomial module
\(D_\tau/D_\tau\partial_\tau\) to the point module
\(\delta_0=D_s/D_s s\), where \(D_u=\mathbb C[u]\langle\partial_u\rangle\).
It would give a quotient of a submodule of \(M\) sending \(m\)
nontrivially to \(\delta_0\). The induced filtration on submodules
and strictness for morphisms
\cite[Proposition~3.19(ii) and Corollary~3.20]{chen-dirks-mustata}
put the image of \(m\) in \(V^\alpha\delta_0\). For \(\alpha>0\)
the point module has zero such piece
\cite[Example~3.18]{chen-dirks-mustata}.
The contradiction proves the lemma.
\end{proof}
